% Day 5's single deck: learning outcomes, the day's plan, and the contents.
\begin{frame}{Learning Outcomes}
    \begin{itemize}
        \setlength\itemsep{0.55em}
        \item Find where bias enters an ML system, audit a model with disaggregated metrics, and explain why the main
              fairness definitions cannot all hold at once.
        \item Explain a model's predictions with permutation importance, SHAP, LIME and Grad-CAM, and know when not to
              trust the explanation.
        \item Document a model with a model card, and place it within SDAIA's AI Ethics Principles and the PDPL.
        \item Take a model from experiment to production: reproducible runs, experiment tracking, and a registry with
              versions and aliases.
        \item Deploy and monitor a model: choose a serving pattern, detect drift before the labels arrive, and release
              behind shadow and canary tests.
        \item Make a model smaller and faster with quantisation, pruning and distillation, and gate the compressed
              version like any other release.
    \end{itemize}
\end{frame}

\begin{frame}{Today: Two Lectures, Two Labs}
    \centering
    \renewcommand{\arraystretch}{1.2}
    {\small
    \begin{tabular}{>{\raggedright\arraybackslash}p{2.0cm}>{\raggedright\arraybackslash}p{9.6cm}}
        \toprule
        \textbf{Block} & \textbf{Content} \\
        \midrule
        Lecture 1 & Part 1: Bias and Fairness \newline Part 2: Interpretability \newline
                    Part 3: Documentation and Governance \\
        Lecture 2 & Part 4: From Experiment to Production \newline Part 5: Lifecycle and Versioning \newline
                    Part 6: Deployment and Monitoring \newline Part 7: Compression and Efficient Inference \\
        Lab 1     & Quantisation: by hand, with ONNX Runtime, then on a YOLO11 detector, live on a webcam \\
        Lab 2     & Pruning and distillation, then ship one through a release gate and an MLflow registry \\
        \bottomrule
    \end{tabular}}
    \vspace{0.5em}
    \begin{block}{Optional labs, on the course site under Day 5}
        A Responsible AI audit (fairness metrics, SHAP and LIME, a model card), and an end-to-end MLOps loop
        (track, tune, deploy, monitor, roll back).
    \end{block}
\end{frame}

\begin{frame}[allowframebreaks]{Table of Contents}
    \tableofcontents
\end{frame}
